\chapter{Teaching a Yoga Class}\label{ch:teaching-class}

The candidate must know how a teacher prepares for a class, what to do at the start and during the class, and the four factors that influence teaching. The notes are read against the official guidebook, an older scheme, as a coverage check. A certified teacher must still review this chapter.

\section{Preparation for a yoga class}

\subsection*{Before the class}

\begin{itemize}
\item Know the practices and the theory you will teach.
\item Plan the order and the time.
\item Set learning objectives with the plan. Objectives, activities and assessment should line up.
\item Arrive early and calm.
\item Check the room.
\end{itemize}

\subsection*{The place}

\begin{itemize}
\item Clean, fresh and well ventilated.
\item Free of noise and bad smell.
\item About the space of a mat for each person, with room to lie down.
\item Mats or blankets on a level floor. Mats are placed in the room or brought by the practitioner, preferably cotton or a poor conductor.
\item Mats and tools are clean and in good order.
\end{itemize}

\subsection*{The student}

\begin{itemize}
\item Empty or light stomach and bowel.
\item Comfortable clothes.
\item A mat of their own is good for hygiene.
\end{itemize}

\subsection*{Opening the class}

\begin{itemize}
\item Greet the students and introduce yourself.
\item Take brief introductions in the first session.
\item Welcome each student and note any personal need.
\item Start in silence, or with a prayer or a short meditation.
\end{itemize}

\subsection*{During the class}

\begin{itemize}
\item Be pleasant and speak clearly.
\item Watch the flow and adjust it from the students' feedback and progress.
\item Admit a mistake with a smile.
\item Keep the class calm.
\item End by thanking the students.
\end{itemize}

The guidebook gives no checklist for other props or tools. Keep to the points above.

\begin{recall}
\begin{itemize}
\item Before means room, plan and self.
\item During means voice, calm and attention.
\item Place: clean, fresh, well ventilated, quiet, a mat of space for each person.
\item Student: light stomach, comfortable clothes.
\item Start with greeting and introductions. End by thanking the students.
\end{itemize}
\end{recall}

\subsection*{Worked examples}

\textbf{How does a teacher prepare for a class?} Use the memory aid: before means room, plan and self; during means voice, calm and attention.
\begin{itemize}
\item Self: know the practices and the theory you will teach; arrive early and calm.
\item Plan: plan the order and the time; set learning objectives so that objectives, activities and assessment line up.
\item Room: clean, fresh, well ventilated, free of noise and bad smell; about a mat of space for each person with room to lie down; mats or blankets on a level floor.
\item During: be pleasant, speak clearly, keep the class calm, and end by thanking the students.
\end{itemize}

\textbf{The first session of a new group.}
\begin{itemize}
\item Arrive early and calm, and check the room and the mats.
\item Greet the students and introduce yourself. Take brief introductions.
\item Welcome each student and note any personal need.
\item Start in silence, or with a prayer or short meditation.
\item As the class goes on, watch the flow and adjust it from the students' feedback and progress.
\item End by thanking the students.
\end{itemize}

\textbf{You make a mistake while cueing.} Admit the mistake with a smile, give the correct cue clearly, and keep the class calm. The teacher is asked to be pleasant and to speak clearly, not to be perfect.

\subsection*{Common questions}

\begin{itemize}
\item \emph{How should a class begin?} Greet the students and introduce yourself, take brief introductions in the first session, and start in silence or with a prayer or short meditation.
\item \emph{Should every student bring a mat?} Mats are placed in the room or brought by the practitioner. A personal mat helps hygiene.
\item \emph{Is this the Board's own teaching guidance?} Not yet. The Board's current teaching text is still to be added.
\end{itemize}

\subsection*{Check yourself}

\begin{enumerate}
\item Fill in the blank: the place for a class should be clean, fresh and \underline{\hspace{2cm}}.
\item What should each person have in the room?
\item How should the teacher end the class?
\item True or false: only the teacher needs to prepare; the student has nothing to prepare.
\item Why is a personal mat good?
\item What does it mean that objectives, activities and assessment should line up?
\item Name two points about the place for a class.
\item What should be checked before the class about the equipment?
\item What does the teacher do about the flow during the class?
\item A new student tells you before class that she has a personal need. What do you do at the start?
\item Halfway through, students look tired and cannot keep up. What does the teacher do?
\end{enumerate}

\subsection*{Answers}

\begin{enumerate}
\item Well ventilated.
\item About the space of a mat, with room to lie down.
\item By thanking the students.
\item False. The student should come with an empty or light stomach and bowel and comfortable clothes. A mat of their own is good for hygiene.
\item It helps hygiene.
\item The learning objectives set should match what is taught and how it is assessed.
\item Any two of: clean and well ventilated, quiet, enough space for each person to lie down.
\item That mats and tools are clean and in good order.
\item Watches it and adjusts it from student feedback and progress.
\item Welcome her and note the personal need, then adjust the class to the needs of each student as it goes on.
\item Watch the flow and adjust it from the students' feedback and progress.
\end{enumerate}

\section{Factors influencing yoga teaching}

Remember four factors: the teacher, the student, the place and the method.

\subsection*{The four factors}

\begin{itemize}
\item \textbf{Teacher:} knowledge, skill in demonstration, communication, attitude and preparation.
\item \textbf{Student:} age, health, level, interest and motivation. People learn in different ways, by seeing, hearing or doing.
\item \textbf{Place (the setting):} space, air, noise, mats, time of day and the size of the group.
\item \textbf{Method:} instruction and demonstration, group or individual teaching, directed practice, discussion. The method must suit the aim.
\end{itemize}

\subsection*{Principles of learning}

\begin{itemize}
\item Interest helps.
\item Learning needs activity and practice.
\item Learning is gradual and individual.
\item Short rests between practices help.
\item Understanding the task before the practice helps the skill.
\end{itemize}

\subsection*{What makes teaching effective}

\begin{itemize}
\item Learners need activity and real experience.
\item The practice itself must be correct.
\item Watching a good demonstration builds skill.
\item Pictures and models help the learning of movement.
\end{itemize}

\subsection*{Methods that work}

\begin{itemize}
\item Combine talk with practice. A lecture works better combined with other methods.
\item Use visuals so that points are remembered.
\item Train the class to respond together. In the response-to-instruction method, the teacher gives precise instructions and all students respond in the same way.
\item Give chances to apply the learning by doing, as in project and directed practice.
\end{itemize}

\subsection*{When you do not know}

The teacher is expected to know more than the students on the points discussed. If unsure, say so and promise to find the answer for the next class.

\begin{recall}
\begin{itemize}
\item Four factors: teacher, student, place, method.
\item People learn by seeing, hearing or doing.
\item The method must suit the aim.
\item Talk works better with practice and visuals.
\item If you do not know, say so and find out.
\end{itemize}
\end{recall}

\subsection*{Worked examples}

\textbf{Name and explain the factors that influence yoga teaching.} Write the four factors first, then one line on each: teacher, student, setting (place) and method, as listed above.

\textbf{A question the teacher cannot answer.} Do not guess. Say honestly that you are not sure. Promise to find the answer and bring it to the next class.

\textbf{Sorting a class problem into a factor.} Ask: teacher, student, place or method? A class too noisy to hear is the place. A student who loses interest is the student. A long talk with no practice is the method. Your own poor preparation is the teacher.

\subsection*{Common questions}

\begin{itemize}
\item \emph{Why does the method have to suit the aim?} The choices work differently, so pick the one that serves what you want the class to learn.
\item \emph{Do pictures and demonstrations really help?} Yes. Watching a good demonstration builds skill, and pictures and models help the learning of movement.
\item \emph{Why do short rests between practices matter?} They are one of the principles that help learning.
\item \emph{Does technology change how people learn?} The guidebook says advances in communication technology have changed teaching systems. It does not say they change how people learn. Treat this as a comment only.
\end{itemize}

\subsection*{Check yourself}

\begin{enumerate}
\item Fill in the blank: the four factors that influence teaching are the teacher, the student, the \underline{\hspace{2cm}} and the method.
\item Name three things about the setting that can help or hinder a class.
\item True or false: people all learn in the same way.
\item Match each to its factor: (a) age and health, (b) skill in demonstration, (c) group or individual teaching, (d) mats and air.
\item Give two principles of learning.
\item Name two points that make teaching effective.
\item What should a teacher do when asked something they do not know?
\item A teacher talks for a long time and the class gets restless. What do you suggest?
\item A teacher wants the whole class to start a practice together on one instruction. Which method fits, and why?
\end{enumerate}

\subsection*{Answers}

\begin{enumerate}
\item Place (the setting).
\item Any three of: space, air, noise, mats, time of day, size of the group.
\item False. People learn in different ways, by seeing, hearing or doing.
\item (a) student, (b) teacher, (c) method, (d) setting.
\item Any two of: interest helps; learning needs activity and practice; learning is gradual and individual; short rests between practices help; understanding the task before the practice helps the skill.
\item Any two of: activity and real experience; correct practice; watching a good demonstration; pictures and models.
\item Say so and commit to find the answer for the next class.
\item A lecture works better combined with other methods. Add practice, use visuals so points are remembered, and let students apply the learning by doing.
\item The response-to-instruction method: the teacher gives precise instructions and all students respond to them in the same way.
\end{enumerate}

\paragraph{Sources.} Official guidebook for Level I (Instructor), Ministry of AYUSH and Quality Council of India, Excel Books, 2016, the teaching practice chapter (older scheme, used as a coverage check); the Board syllabus for Level 1, Yoga Protocol Instructor. In own words.
